\section{Consistency lemmas of the definitions}
\label{supp:lemmas}

The following lemmas are proved by projecting the fields of the definitions D3--D6 of the main paper
or by contradiction with one of them. Their use is that later certificates cannot silently omit a
detector, a gate, a transition, an audit or a source check.

\begin{lemma}[source exclusion]\label{lem:d3}
A record whose source tag is fallback, independent-pair witness, simulation trace, ablation record,
unknown or contradictory is not carried, whatever its flags. A record reached through a step outside
the declared scope is not carried. An instrument whose record list is not declared complete is not a
carried instrument, even if every listed record is carried.
\end{lemma}

\begin{proof}
Each clause contradicts one conjunct of \Cref{def:d3}.
\end{proof}

\begin{lemma}[repair-step consistency]\label{lem:d4}
A lawful repair step entails detection, gate permission, a kernel transition and re-audit; a
candidate step lacking any one of the four is not lawful. No repair move has sort P6.
\end{lemma}

\begin{lemma}[closed-loop consistency]\label{lem:d5}
In a closed-loop history every listed repair entry is carried. One listed entry with source tag
fallback, unknown or contradictory, or with a false generation or scope flag, defeats closed-loop
scope. A fallback entry is not carried even when a different entry with the same payload is.
\end{lemma}

\begin{lemma}[probe-move consistency]\label{lem:d6}
A lawful allocation does not change the support; a lawful acquisition is strict; a lawful retirement
removes a probe that was active; and an operator-saturated probe is not a strict acquisition.
\end{lemma}
